% capitoli/03_metodologia.tex - Capitolo 3: Metodologia

\chapter{Metodologia}
\label{ch:metodologia}

\section{Estrazione e Preprocessing dei Dati}
Il punto di partenza dell'esperimento è costituito da log ospedalieri estratti in formato XES. L'obiettivo primario di questa fase iniziale è strutturare e formattare tali log non strutturati per renderli digeribili all'architettura neurale sottostante il Transformer. 

\subsection{Analisi Esplorativa del Dataset (EDA)}
Per acquisire contezza della natura e del volume dei log estratti pre-elaborazione, è stata condotta un'Analisi Esplorativa dei Dati (EDA) sviluppando una pipeline automatizzata in Python (basata principalmente sulla libreria \texttt{pandas} per l'efficiente gestione e \textit{chunking} massivo delle estrazioni). Il dataset grezzo comprende 7.393 casi clinici unici (singoli pazienti ospitati), all'interno dei quali sono registrati 88.115 eventi granulari totali non etichettati.

Si rileva contestualmente l'ingente sbilanciamento delle classi, cardine della sfida analitica affrontata in questo progetto: su 7.393 pazienti totali, ben 6.584 esibiscono una degenza classificabile come normale (classe 0), e solamente 809 costituiscono campioni d'anomalia con degenza critica $\ge 20$ giorni (classe 1).

\begin{figure}[htbp]
    \centering
    \includegraphics[width=0.8\textwidth]{"img/actions_per_case_distribution.png"}
    \caption{Distribuzione statistica globale del numero di azioni processuali registrate per singolo caso clinico. Il grafico si caratterizza per una spiccata asimmetria (\textit{right-skewed}), evidenziando come una concentrazione primaria di interazioni ospedaliere rapide all'ingresso si contrapponga ad una ``lunga coda'' indicante pazienti complessi trattenuti in cure iterative stratificate.}
    \label{fig:actions_distribution}
\end{figure}
\FloatBarrier

Come illustrato in Figura \ref{fig:actions_distribution}, il numero di azioni per paziente presenta una forte asimmetria destra (\textit{right-skewness}): la media è di 12.92 azioni, mentre la mediana scende a 9.00, segnalando che la distribuzione è tirata verso l'alto da una ``lunga coda'' (\textit{long tail}) di pazienti complessi con decine o centinaia di re-interazioni.

\begin{figure}[htbp]
    \centering
    \includegraphics[width=0.8\textwidth]{"img/class_quantity_boxplot.png"}
    \caption{Boxplot di confronto diretto tra il numero d'interazioni erogate ai pazienti di classe 0 (degenza limitata) e di classe 1 (degenza prolungata). Risalta ad occhio nudo la traslazione dei quartili a sfavore della classe anomala, gravata intrinsecamente da percorsi ben più caotici, iterativi ed onerosi quantitativamente per la clinica.}
    \label{fig:class_quantity_boxplot}
\end{figure}
\FloatBarrier

Tale diversità emerge chiaramente analizzando le durate per classe (Figura \ref{fig:class_quantity_boxplot}): appare evidente il divario numerico, con la classe 1 (long-LOS) caratterizzata da percorsi mediani molto più lunghi e stratificati rispetto ai pazienti con degenti brevi. Questa evidenza giustifica l'adozione di un'architettura avanzata come BERT. A differenza di una stringa Markoviana classica che frammenterebbe i calcoli su sequenze enormi, l'architettura basata su \textit{Self-Attention} garantisce un'elaborazione in parallelo per mappare le dipendenze logiche tra decine di step ospedalieri.

\begin{figure}[htbp]
    \centering
    \includegraphics[width=\textwidth]{"img/class_action_prevalence_diff_top.png"}
    \caption{Panoramica estrapolata dei \textit{delta} di prevalenza percentuale (sbilanciamento) per le singole azioni ospedaliere fortemente correlate in modo endogeno alla classe d'emergenza logistica 1.}
    \label{fig:class_action_prevalence}
\end{figure}
\FloatBarrier

Un'ulteriore esplorazione evidenzia la distribuzione asimmetrica delle frequenze delle singole procedure (Figura \ref{fig:class_action_prevalence}). Risultano particolarmente pronunciati gap associati ad alta probabilità critica per eventi clinici come la ``Visita Infettivologica'', la ``Radiografia Torace a Letto'' o il ``Posizionamento Catetere Venoso Centrale''. Pur trattandosi di statistiche esplorative prive di valore causale formale, tali osservazioni anticipano i pattern che il modello imparerà a pesare durante il training e che verranno interpretati tramite XAI nel Capitolo \ref{ch:risultati}.

\subsection{Approcci e Formattazione degli Embedding}
Poiché il task ingegneristico mira strettamente ad accoppiare eventi clinici elaborati a corretto formato ed interfacciandoli poi col \textit{Learning} del Transformer, è stata condotta una meticolosa opera preliminare testando tre filosofie distinte di \textit{embedding}:

\begin{enumerate}
    \item \textbf{Azioni cliniche + Binning Temporale:} In questo approccio, il preprocessing concatena coppie formate dall'azione clinica e da un token temporale speciale. Il timing continuo tra le azioni viene discretizzato (\textit{binning}) in categorie predefinite. Un esempio di traccia processata è il seguente:
\begin{minted}[fontsize=\small]{json}
{"case_id": "5870415", 
 "tokens": ["Complete abdomen ultrasound", "_TIME_DELTA_START_", 
            "Computerized tomography of the skull", 
            "_TIME_DELTA_START_", "Complete abdomen ultrasound", 
            "_TIME_DELTA_LONG_", "Computerized tomography of the skull", 
            "_TIME_DELTA_MEDIUM_"], 
 "label_id": 0
}
\end{minted}

    La scelta delle soglie di discretizzazione non è arbitraria ma \textit{data-driven}: le categorie sono state definite analizzando la distribuzione empirica dei delta temporali nel dataset, riportata in Tabella~\ref{tab:time_delta_distribution}. La distribuzione rivela una struttura bimodale con la grande maggioranza degli intervalli concentrata entro la finestra dei 7 giorni (97.0\%), e una coda esigua di eventi molto dilazionati. Le soglie sono state quindi fissate sui breakpoint clinici naturali --- un'ora, un giorno, sette giorni e trenta giorni --- che separano fasi di cura operativamente distinte: accertamenti rapidi, degenza standard, ricovero prolungato e cronicità.

    \begin{table}[htbp]
        \centering
        \resizebox{\textwidth}{!}{%
        \begin{tabular}{|l|c|r|r|}
            \hline
            \textbf{Token} & \textbf{Intervallo} & \textbf{Conteggio} & \textbf{Freq.} \\
            \hline
            \texttt{\_TIME\_DELTA\_START\_}      & $\Delta = 0$ s                         & 5\,842  &  6.9\% \\
            \texttt{\_TIME\_DELTA\_IMMEDIATE\_}  & $0 < \Delta \le 1\text{h}$             & 19\,456 & 23.1\% \\
            \texttt{\_TIME\_DELTA\_SHORT\_}      & $1\text{h} < \Delta \le 1\text{gg}$   & 29\,887 & 35.4\% \\
            \texttt{\_TIME\_DELTA\_MEDIUM\_}     & $1\text{gg} < \Delta \le 7\text{gg}$  & 26\,649 & 31.6\% \\
            \texttt{\_TIME\_DELTA\_LONG\_}       & $7\text{gg} < \Delta \le 30\text{gg}$ &  2\,296 &  2.7\% \\
            \texttt{\_TIME\_DELTA\_VERY\_LONG\_} & $\Delta > 30\text{gg}$                 &    233  &  0.3\% \\
            \hline
            \multicolumn{2}{|l|}{\textbf{Totale delta calcolati}} & 84\,363 & 100\% \\
            \hline
        \end{tabular}%
        }
        \caption{Distribuzione empirica dei delta temporali nel dataset e corrispondente token di discretizzazione assegnato. Il totale è inferiore agli 88.115 eventi grezzi poiché il primo evento di ciascuna traccia non possiede un evento precedente da cui calcolare il delta.}
        \label{tab:time_delta_distribution}
    \end{table}
    \FloatBarrier

    La funzione \texttt{categorize\_time\_delta\_vectorized}, implementata con \texttt{numpy} per un processing completamente vettorializzato sull'intero dataset, traduce i delta continui nei token speciali secondo questa logica:

\begin{minted}[fontsize=\small]{python}
def categorize_time_delta_vectorized(deltas):
    conditions = [
        deltas == 0,                               # Delta nullo: eventi simultanei
        (deltas > 0) & (deltas <= 3600),           # 0-1 ora:   Immediato
        (deltas > 3600) & (deltas <= 86400),       # 1h-1gg:    Breve
        (deltas > 86400) & (deltas <= 604800),     # 1-7gg:     Medio
        (deltas > 604800) & (deltas <= 2592000),   # 7-30gg:    Lungo
    ]
    choices = [
        '_TIME_DELTA_START_',       # delta = 0
        '_TIME_DELTA_IMMEDIATE_',   # 0 < delta <= 1h
        '_TIME_DELTA_SHORT_',       # 1h < delta <= 1gg
        '_TIME_DELTA_MEDIUM_',      # 1gg < delta <= 7gg
        '_TIME_DELTA_LONG_',        # 7gg < delta <= 30gg
    ]
    return np.select(conditions, choices, default='_TIME_DELTA_VERY_LONG_')
\end{minted}

    \item \textbf{Array concatenati (Azioni + Distanze numeriche):} Questa seconda forma di preprocessing costruisce l'embedding affiancando due array distinti: uno contenente esclusivamente la sequenza testuale delle azioni cliniche e un array parallelo numerico (\textit{float}) esprimente la distanza in secondi dall'evento precedente. Un esempio:
\begin{minted}[fontsize=\small]{json}
{"case_id": "5870415", 
 "tokens": ["Complete abdomen ultrasound", "Computerized tomography of the skull", 
            "Complete abdomen ultrasound", "Computerized tomography of the skull"], 
 "time_deltas": [0.0, 0.0, 728160.0, 178980.0], 
 "label_id": 0
}
\end{minted}

    \item \textbf{Storytelling in linguaggio naturale:} L'approccio finale, che costituisce il core metodologico del progetto TEXLOS \cite{texlos2026} e che è stato preferito ai precedenti, si rifà al pionieristico lavoro di traduzione semantica dei log del progetto LEGOLAS \cite{PASQUADIBISCEGLIE2025128224}. Anziché alimentare il classificatore con una combinazione mista di vettori numerici e codici, gli eventi e i vincoli cronologici logici della traccia sono tradotti letteralmente in frasi inglesi strutturate grammaticalmente, producendo interi paragrafi coerenti narranti il percorso diagnostico/terapeutico di ciascun paziente.
\begin{minted}[fontsize=\small]{json}
{"case_id": "5870415", 
 "story": "The patient, 87 years old, underwent a series of examinations and 
 interventions during hospitalization.
After 0 seconds, the following examinations were performed simultaneously: 
Complete abdomen ultrasound and Computerized tomography of the skull.
After 728,160 seconds, Complete abdomen ultrasound was performed.
After 178,980 seconds, Computerized tomography of the skull was performed.",
"label_id": 0
}
\end{minted}
\end{enumerate}

\subsection{Dettagli Implementativi dell'Algoritmo di Storytelling}
Per colmare la distanza che intercorre tra i log estratti in formato grezzo e i testi discorsivi forniti in addestramento, è stata sviluppata un'infrastruttura di \textit{parsing} ed iniezione in linguaggio Python. Dal punto di vista informatico, il processo si divide in due fasi computazionali: la strutturazione in memoria del file XES e la generazione cronologica della storia narrata. 

Nella prima fase entra in gioco il modulo \texttt{xes\_parser.py}. Affrontando l'ingente volume di eventi clinici asincroni tipici di un'analisi di Process Mining, l'algoritmo si destreggia tra la libreria per l'ingestione dei log (\texttt{pm4py}) e le più ottimali e leggere architetture dati di \texttt{pandas}. Quest'ultimo processa massivamente il dataset suddividendo mediante \textit{groupby} e parallelismi nativi le righe afferenti ai medesimi pazienti; si procede poi convertendo in formato standard i \textit{timestamp} (le date d'esecuzione dell'atto medico), le stringhe degli attori coinvolti ed instradando così ogni frammento cronologico entro uno strato informativo solido (oggetto \textit{PatientTrace}).  

La fase autoriale vera e propria è governata dalla classe \texttt{StoryGenerator} (all'interno dello script \texttt{story\_generator.py}). Le specifiche di traduzione che intercettano queste \textit{trace} vertono su tre stringenti requisiti ingegneristici:
\begin{itemize}
    \item \textbf{Gestione del Tempo (Time Deltas):} L'evoluzione clinica viene misurata estraendo rigorosamente il \textit{delta temporale in secondi} che separa l'azione attuale dalla precedente operazione registrata. Questo gap continuo viene iniettato nel template e calcolato esplicitamente mediante un blando \textit{parsing} aritmetico (\texttt{current\_time - prev\_time}), oggettivando la scansione temporale. Nessuna approssimazione logica interviene nel testo finale; i secondi marcano senza preconcetti le tempistiche di risposta dei triage o la dilatazione cronica della degenza (si passa da "0 secondi" per l'esordio iniziale, a svariate centinaia di migliaia di secondi per tac e risonanze d'esito fatte la settimana d'uscita).
    \item \textbf{Gestione degli Eventi Simultanei:} Non è inusuale, nei processi su scala giornaliera densi come in ospedale, rilevare molteplici interazioni registrate col medesimo identico \textit{timestamp}. Un esempio classico si verifica con batterie di check-up del sangue incodificate alla stessa ora informatica dal personale. Per evitare grave ridondanza lessicale che sovraccaricherebbe indebitamente l'asse d'attenzione di BERT con frasi ripetute, il generatore intercetta internamente gli aggregati (tramite la funzione \texttt{\_group\_simultaneous\_events}), fondendoli all'interno di specifici \textit{bullet points} o paragrafi discorsivi compatti nel seguente formato: \textit{"After $X$ seconds, the following examinations were performed simultaneously: Action A, and Action B"}.
    \item \textbf{Indipendenza dal Protocollo di Triage:} Il generatore proposto si differenzia strutturalmente dal progetto LEGOLAS, il quale era progettato per interpolare un insieme fisso di variabili fisiologiche standard da Pronto Soccorso (frequenza respiratoria, temperatura corporea, scala del dolore, ecc.), compilando un template narrativo predefinito. Nel contesto di TEXLOS, i pazienti attraversano reparti eterogenei con procedure cliniche non standardizzate, prive di una nomenclatura universale codificata. Per gestire questa variabilità, il generatore adotta dizionari interni flessibili e \textbf{regole di template matching dinamiche}, agnostiche rispetto alla tipologia di azione clinica, che mappano la tassonomia locale dell'ospedale ai costrutti discorsivi del Transformer senza vincolare l'algoritmo a un sottoinsieme predefinito di sensori o misurazioni.
\end{itemize}

\newpage
A scopo esemplificativo, il seguente frammento di pseudo-codice descrive la routine adoperata dal software in \texttt{story\_generator.py} per snellire testualmente (\textit{trimming} sintattico) il sovraffolamento degli stadi simultanei ed iniettare i \textit{time-delta} per una traccia degenziale prima dell'invio in NLP ad estrapolazione avvenuta:

\begin{minted}[fontsize=\small]{python}
def create_narrative_from_trace(trace_events):
    sorted_events = sort_by_timestamp(trace_events)
    grouped_events = group_simultaneous_events(sorted_events)
    
    narrative_paragraphs = []
    prev_time = grouped_events[0].timestamp
    
    for group in grouped_events:
        time_elapsed = calculate_seconds_diff(prev_time, group.timestamp)
        
        if len(group) == 1:
            # Singolo evento testuale formattato
            activity_desc = apply_medical_template(group[0].activity)
            text = f"After {time_elapsed} seconds, {activity_desc} was performed."
        else:
            # Piu' azioni collassate per alleggerimento dell'Attention
            activities = [apply_medical_template(e.activity) for e in group]
            concat_acts = join_with_commas_and(activities)
            text = f"After {time_elapsed} seconds, the following examinations " \
                   f"were performed simultaneously: {concat_acts}."
                   
        narrative_paragraphs.append(text)
        prev_time = group.timestamp # Sposta l'ancora
        
    return "\n\n".join(narrative_paragraphs)

\end{minted}

\section{Architettura del Modello e Limitazioni dei BERT Clinici}
Alla base del sistema di previsione giace \texttt{bert-base-uncased}, un Large Language Model reso disponibile pubblicamente sulla piattaforma HuggingFace. Questo Transformer, limitato alla componente Encoder, viene sottoposto a Fine Tuning aggiungendo al termine dell'architettura standard una testa di classificazione (\textit{classification head}) volta a predire la natura binaria del problema: degenza normale (sotto i 20 giorni) contro degenza critica (uguale o superiore a 20 giorni).

Il primo istinto sperimentale è stato chiaramente quello di testare modelli \textit{Domain-Specific}, confrontando le prestazioni di varianti pre-addestrate unicamente su corpus clinici -- quali \texttt{biobert}, \texttt{pubmedBert}, \texttt{blue-bert} e \texttt{clinical-bert} -- con quelle di \texttt{bert-base-uncased}. L'ipotesi sottostante era che un vocabolario medico radicato \textit{ab initio} avrebbe accelerato la decodifica delle diagnosi. L'evidenza empirica ha sorprendentemente sancito prestazioni sovrapponibili tra i vari modelli e, anzi, una \textit{Balanced Accuracy} leggermente inferiore per i modelli clinici rispetto al canonico modello generalista. 

Tale anomalia si spiega analizzando l'altissima specificità terminologica del dataset della tesi: i modelli addestrati su dataset clinici (come MIMIC) sono spesso altamente specifici e ottimizzati per una formale nomenclatura standard inglese. Nel nostro caso, il nomenclatore locale dell'ospedale era ricco di acronimi, contrazioni proprietarie e formulazioni non standardizzate, vanificando almeno parzialmente il vantaggio della pregressa conoscenza medica dei pesi. Questo disallineamento è critico soprattutto a livello di \textit{tokenizer}: i tokenizer dei vari modelli clinici, agendo su un vocabolario molto più verticale, frammentavano le parole spurie del dataset della tesi perdendo efficacia. Di contro, il tokenizer base di \texttt{bert-base-uncased} si è dimostrato molto più elastico, resiliente e capace di preservare sufficiente compattezza semantica pur al cospetto di un dominio estraneo.

\section{Giustificazione Architetturale: Analisi del Trade-off Computazionale}
In fase di definizione dell'architettura neurale, è stata condotta un'indagine esplorativa scalando il modello base verso la sua variante maggiorata \texttt{bert-large-uncased} (circa 340 milioni di parametri contro i 110 milioni originari). L'intento era tracciare la scalabilità dell'accuratezza su una rete più stratificata.

Un test preliminare su una singola esecuzione ha restituito metriche di rilievo sulla classe critica (LOS $\ge 20$ giorni), registrando una Precision del 73.65\% e una Recall dell'84.36\%. Tuttavia, pur evidenziando un lieve incremento prestazionale grezzo rispetto alla versione base, l'esperimento ha fatto emergere severi limiti di scalabilità e di \textit{data plateauing}. Da un lato, l'aumento dimensionale del modello non ha prodotto un salto di qualità proporzionale al costo computazionale, suggerendo che un bacino di circa 7.000 tracce cliniche non sia sufficiente a saturare la capacità di astrazione di un'architettura "Large", portando di fatto la rete in una fase di plateau informativo.

Dall'altro lato, i vincoli hardware hanno reso la sperimentazione ingegneristicamente insostenibile. La saturazione della VRAM ha imposto una drastica riduzione della \textit{batch size}, precipitata dai 22 campioni gestibili dal modello base ad appena 4 campioni per il modello large. Questo crollo ha dilatato esponenzialmente i tempi di addestramento, rendendo materialmente inapplicabile la \textit{Stratified 5-Fold Cross Validation}, requisito imprescindibile delineato in questa tesi per garantire la robustezza statistica del classificatore dinanzi a uno sbilanciamento di classi così severo.

Di conseguenza, il modello \texttt{bert-base-uncased} è stato eletto come "ottimo paretiano" del progetto, offrendo il compromesso ideale tra capacità predittiva, sostenibilità hardware e garanzia di validazione metodologica rigorosa.

\section{Focal Loss e Sbilanciamento delle Classi}
Il dataset presenta uno dei paradossi strutturali più frequenti della Process Analytics in sanità: lo \textit{sbilanciamento estremo delle classi}. Di fatto, le casistiche in cui il paziente sfiora o sfora una degenza continua di venti giorni (classe 1 - Critico) costituiscono una netta e fisiologica minoranza anomala nei log degli eventi rispetto alla travolgente maggioranza dei passaggi di routine o dimissioni rapide (classe 0 - Normale).

Una prima fase di addestramento gestita dalla standard \textit{Binary Cross Entropy} (BCE) ha confermato l'incapacità del modello di apprendere in modo bilanciato: il classificatore tendeva a privilegiare la classe dominante, producendo di fatto una predizione prossima al \textit{ZeroR} e penalizzando la sensibilità sui casi critici — esattamente quelli di primario interesse per le amministrazioni ospedaliere.

Per correggere questo bias strutturale e ricalibrare il contributo al gradiente in modo differenziale, è stata adottata in fase di fine-tuning la funzione di costo denominata \textbf{Focal Loss} \cite{lin2017focal}. Questa funzione mutua la sua origine nel dominio della \textit{Computer Vision}, dove nacque per correggere il massiccio sbilanciamento tra uno sconfinato spazio di \textit{background} e piccoli oggetti da identificare. Applicata al processing di tracce cliniche, la Focal Loss riduce il peso computazionale (\textit{loss contribution}) degli esempi facili e sovrabbondanti, in favore dei campioni mal classificati.

\subsection{Formulazione Matematica della Focal Loss}
Il punto di partenza è la funzione di costo \textit{Cross Entropy} (CE) per la classificazione binaria. Definito $y \in \{\pm 1\}$ come la classe ground-truth e $p \in [0, 1]$ come la probabilità stimata dal modello per la classe $y=1$, la formulazione classica è:

\begin{equation}
CE(p, y) = \begin{cases} -\log(p) & \text{se } y=1 \\ -\log(1-p) & \text{altrimenti.} \end{cases}
\end{equation}

Per convenienza notazionale, definendo $p_t$:
\begin{equation}
p_t = \begin{cases} p & \text{se } y=1 \\ 1-p & \text{altrimenti,} \end{cases}
\end{equation}
la CE si può riscrivere compattamente come $CE(p_t) = -\log(p_t)$.

Un metodo comune per mitigare lo sbilanciamento puro introduce un fattore di peso $\alpha \in [0, 1]$ per la classe 1 e $1-\alpha$ per la classe -1. Definendo $\alpha_t$ in modo analogo a $p_t$, si ottiene la \textit{$\alpha$-balanced Cross Entropy}:
\begin{equation}
CE(p_t) = -\alpha_t \log(p_t)
\end{equation}

Tuttavia, pur bilanciando l'importanza numerica tra positivi e negativi, questa formulazione non differenzia tra esempi \textit{facili} (\textit{easy negatives}) ed esempi \textit{difficili} (\textit{hard examples}). Per abbattere il peso sproporzionato degli innumerevoli esempi facili (i pazienti con degenze standard facilmente prevedibili), Lin et al. introducono un fattore modulante $(1-p_t)^\gamma$, dove $\gamma \ge 0$ è un iperparametro noto come \textit{parametro di focalizzazione} (\textit{focusing parameter}). La formulazione completa della \textbf{Focal Loss} (nella sua variante $\alpha$-balanced pratica) diviene quindi:

\begin{equation}
FL(p_t) = -\alpha_t (1-p_t)^\gamma \log(p_t)
\end{equation}

Questa funzione gode di proprietà analitiche fondamentali per il nostro task:
\begin{itemize}
    \item Se un esempio ospedaliero è mal classificato o raro (quindi $p_t$ è piccolo), il fattore modulante $(1-p_t)^\gamma$ tende a 1 e la \textit{loss contribution} rimane inalterata, garantendo l'apprendimento.
    \item Se un esempio è facile e ben classificato ($p_t \rightarrow 1$), il fattore modulante tende a 0, penalizzando (\textit{down-weighting}) fortemente l'impatto sul gradiente globale.
    \item Il parametro $\gamma$ regola il tasso di questa penalizzazione (addolcendo l'iper-dominanza della classe 0 nel nostro dataset), esonerando l'analista dalla necessità di applicare euristiche forzate e distruttive a valle, come un aggressivo \textit{undersampling} che distruggerebbe preziosi pattern temporali globali.
\end{itemize}

L'effetto della parametrizzazione $\gamma$ si evince analizzando visivamente l'andamento della funzione per vari scenari, confrontando la Focal Loss con la Cross Entropy standard (Figura \ref{fig:focal_loss_comparison}).

\begin{figure}[htbp]
    \centering
    \includegraphics[width=0.8\textwidth]{"img/focal loss fig 1.png"}
    \caption{Confronto tra la Cross Entropy (CE) e diverse configurazioni della Focal Loss. Sull'asse delle ascisse è riportata la probabilità effettiva stimata dal modello, $p_t \in [0, 1]$, mentre sull'asse delle ordinate figura il valore della funzione di costo. Si noti come, al crescere dell'iperparametro di focalizzazione $\gamma$ (es. $\gamma=2, 5$), la \textit{loss} addebitata agli esempi facili e ben classificati ($p_t \rightarrow 1$) subisca un drastico schiacciamento rispetto al decadimento più lento della CE. Questo espediente permette all'osservazione di focalizzarsi dinamicamente sugli esempi più critici ed errati.}
    \label{fig:focal_loss_comparison}
\end{figure}
\FloatBarrier

Una volta garantita la capacità discriminatoria del modello su classi sbilanciate, si procede all'analisi della sua interpretabilità tramite le tecniche descritte nel capitolo successivo.

\section{Setup Sperimentale e Hyperparameter Tuning}
L'intera pipeline di esperimenti e addestramento supervisionato è stata materialmente elaborata e portata a convergenza su un setup hardware performante, incentrato su una workstation munita di CPU AMD Ryzen 9 5900X, 64 GB di RAM e GPU NVIDIA GeForce RTX 4090 (consentendo l'aumento efficiente delle batch size e l'addestramento distribuito del grafo tensoriale Transformer).

Per preservare l'integrità e il realismo della diagnostica probabilistica al cospetto dello sbilanciamento delle classi, la valutazione accademica si basa unicamente su iterazioni cicliche di \textit{Stratified K-Fold Cross Validation}. Questo garantisce che ogni singola split prodotta preservi identiche proporzioni di pazienti anomali/ordinari rispetto al database globale, smaltendo bias d'addestramento.

L'impostazione degli iperparametri del modello neurale è sfociata da un processo combinato: 
\begin{itemize}
    \item L'esplorazione automatica bayesiana condotta massivamente dal framework \texttt{Optuna}.
    \item Un minuzioso \textit{tuning manuale} a posteriori. Affinché i risultati non scivolassero in asimmetrie dimensionali non volute, le alterazioni manuali degli intervalli innescati da Optuna si sono preoccupate di preservare costanti le "proporzioni architetturali" auree del transformer, quali i rapporti intrinseci tra hidden layers, attention heads e scaling rate dello stack.
\end{itemize}

\subsection{Convergenza del Modello e Early Stopping}
A seguito del setup definito, l'addestramento della configurazione finale \texttt{bert-base-uncased} ha evidenziato un comportamento stabile. Come illustrato nella Figura \ref{fig:loss_curve}, l'andamento mostra una discesa costante della \textit{training loss}, mentre la \textit{validation loss} raggiunge il suo minimo all'epoca 10. Nelle epoche successive la curva di validazione inizia a divergere, indicando l'innesco dell'\textit{overfitting}. Grazie alla tecnica dell'\textit{early stopping} (impostata con un $\Delta$ minimo di $1\times 10^{-3}$ su una \textit{patience} di 5 epoche), il training si è interrotto all'epoca 15; i pesi del modello sono stati quindi ripristinati all'ottimo dell'epoca 10, garantendo la massima capacità di generalizzazione sui dati non visti per la successiva fase di test.

% STORIA LOSS
\begin{figure}[htbp]
    \centering
    \includegraphics[width=0.8\textwidth]{img/train_val_loss.png}
    \caption{Curva di apprendimento (Training vs Validation Loss) del modello definitivo durante il processo di \textit{fine-tuning}.}
    \label{fig:loss_curve}
\end{figure}
\FloatBarrier

\section{Indagine Esplorativa: Data Augmentation con LLM}
Nel tentativo di robustire l'astrazione di BERT sulla classe minoritaria (long-LOS), si è tentato un originale esperimento di estensione del dataset (\textit{Data Augmentation}) avvalendosi dell'intelligenza generativa degli attuali LLM generativi (decoder-only).

La pipeline esplorativa prevedeva l'impiego locale di modelli generativi allo stato dell'arte, nello specifico \texttt{phi4:14b} e \texttt{nous-hermes2:34b}, eseguiti tramite l'ambiente open-source \texttt{Ollama}. Per inibire le allucinazioni e circoscrivere l'effusione creativa dell'esposizione testuale, la \textit{Temperature} è stata fissata al minimo rigido ($T=0.0$). 

Il layer di codice Python istruiva severamente le reti a limitarsi ad un mero lavoro di riformattazione narrativa, passando iterativamente le coordinate base della traccia log:

\begin{minted}[fontsize=\scriptsize]{python}
prompt = f"""You are an expert physician who writes discursive clinical
narratives in English for BERT/NLP analysis.

IMPORTANT - MANDATORY MEDICAL TERMINOLOGY:
You must use EXACTLY the medical terminology specified in the clinical vocabulary.

PATIENT DATA:
- Information: {patient_desc}

MEDICAL EVENTS SEQUENCE (already standardized):
{events_text}

SPECIFIC RULES:
1. USE EXACTLY the provided activity names (already standardized)
2. Introduction with age and medical context
3. Separate paragraph for each temporal group
4. Format numbers: "13,380 seconds" (with commas for thousands)
5. Concurrent events: "the following examinations were performed simultaneously:"
6. Connectors: "The [activity] was performed", "After X seconds [activity] was executed"
7. Paragraphs separated by blank line
8. Keep medical terminology IDENTICAL to provided terminology
9. Always conclude with "At the end of the clinical pathway, the medical events were concluded."
10. NO translations or paraphrases of medical activity names
11. All text must be in ENGLISH"""
\end{minted}

Nonostante l'imbragatura direzionale e sintattica imposta a livello di framework, l'ingestione e la rielaborazione da parte delle reti ha portato a delle devianze gravose. A titolo di dimostrazione, viene qui incollata la narrativa sintetica prodotta dal modello \texttt{phi4:14b} a fronte del caso ospedaliero ID "5870415" (contenente originariamente 4 atti nosocomiali primari ed oggettivi estratti dai registri):

\begin{minted}[fontsize=\scriptsize]{json}
{
  "case_id": "5870415",
  "story_text": "The patient, 87 years old, underwent a series of examinations...
The Complete abdomen ultrasound was performed as the initial examination to 
assess abdominal structures for any abnormalities or pathologies that might 
require further investigation... This non-invasive imaging technique provided 
valuable insights... After 0 seconds, the following examinations were performed 
simultaneously: Computerized tomography of the skull and Complete abdomen ultrasound. The Coputerized tomography of 
the skull was conducted to evaluate any potential intracranial issues such as 
hemorrhages, tumors, or other abnormalities... After 728,160 seconds 
(approximately 8 days and 7 hours), another Complete abdomen ultrasound was 
executed. This follow-up examination aimed to monitor any changes...
At the end of the clinical pathway, the medical events were concluded."
}
\end{minted}

L'esempio documenta i limiti intrinseci nella generazione del testo impiegata per questo scopo predittivo. Nonostante le rigide indicazioni tecniche nel prompt testuale, il modello \texttt{phi4:14b} ha generato un'ampia \textbf{allucinazione semantica}, arricchendo indebitamente l'output testuale: il LLM ha aggiunto descrizioni procedurali e teoriche arbitrarie sul motivo degli esami (\textit{"to evaluate any potential intracranial issues such as hemorrhages"}), che risultano rumore puro ai fini di classificazione ospedaliera.

Test successivi svolti con un LLM dotato di un maggior numero di parametri (\texttt{nous-hermes2:34b}) hanno prodotto storie più sintetiche e rispettose delle regole sintattiche del prompt, immettendo tuttavia alterazioni concettuali sottili e altrettanto problematiche:

\begin{minted}[fontsize=\scriptsize]{json}
{
  "case_id": "5870415",
  "story_text": "The patient, an 87-year-old individual, underwent a series of 
examinations... The first event was a complete abdomen ultrasound... 
Simultaneously, a Computerized tomography scan of the skull without contrast medium was executed. 
After 728160 seconds, another complete abdomen ultrasound was conducted. 
Following this, a Computerized tomography scan of the skull with contrast medium was carried out..."
}
\end{minted}

In questo secondo test, il modello linguistico ha dedotto arbitrariamente dettagli non presenti nel database clinico, specificando la presenza o l'assenza del contrasto in una diagnostica di routine (\textit{``with contrast medium / without contrast medium''}). Benché semanticamente pertinenti in un caso cardio-vascolare, queste integrazioni scostano la traccia artificiale dall'estrazione originale, alterando il vocabolario disponibile a valle e compromettendo la confrontabilità con le tracce reali.

La convalida empirica e i successivi processi di \textit{fine-tuning} incrociato hanno sancito il fallimento dell'ipotesi predetta. L'avvalersi dello \textit{Storytelling} sintetico arricchito da tali artefatti informativi ha comportato un degrado strutturale delle metriche computazionali. La rumorosità lessicale iniettata dai Large Language Models confonde l'attenzionalità temporale del Transformer, peggiorandone la capacità selettiva e compromettendo di fatto i valori di \textit{Balanced Accuracy}. Per via di queste problematiche, l'adozione testuale definitiva in questa tesi usa esclusivamente script di parsing lineari (\texttt{story\_generator.py}), assicurando una perfetta aderenza matematica ed enumerabile al tracciato loggato originario.
